% ============================================================
% PATH PLANNING  (Requirement 4b)
% ============================================================
\section{Path Planning}
\label{sec:planning}

The planner runs once, before the vehicle moves, in five stages: inflate
the walls, route every pair of points with A* to obtain driving
distances, order the waypoints, route each leg in that order, and smooth
each leg with particle swarm optimisation (PSO). As the task requires,
it receives only the wall map and has no knowledge of the obstacles. The
obstacle list is read only when the random waypoints are drawn, to keep
them clear of obstacles; Fig.~\ref{fig:replan} shows a planned path
that runs inside the keep-out of an obstacle it did not know. Obstacles are
handled during the mission (Section~\ref{sec:avoidance}).

\subsection{Map Inflation}
Growing every wall by at least $d_{\text{clear}}$ lets the vehicle be planned as
a point in configuration space \cite{lozanoperez1983, lavalle2006}. The
map is an occupancy grid with $0.1$~m cells, and the planner searches
over the cell corners. Each wall cell is therefore first extended to the
four corners it touches, so that distances from those corners are true
distances from the wall faces. The walls are grown by $0.60$~m (six cells), not
$0.40$~m. The extra $0.20$~m is a tracking margin: the path is followed
from a noisy state estimate, and a path planned at exactly $0.40$~m
would breach the clearance with any tracking error.

\subsection{Working Principle}
\textbf{Routing.} Each leg is routed with A* \cite{hart1968} on the
inflated grid. A* expands the grid point with the lowest
\begin{equation}
    f(n) = g(n) + h(n),
    \label{eq:astar}
\end{equation}
where $g(n)$ is the path length from the start to point $n$ and $h(n)$
is the straight-line distance from $n$ to the goal. Points are
eight-connected. Because $h$ never overestimates the remaining distance,
the first path to reach the goal is the shortest one on the grid.

\textbf{Smoothing.} An A* path is made of straight runs joined at
$45^{\circ}$ corners and hugs the inflated walls (Fig.~\ref{fig:plan});
a steered vehicle cannot follow it. Each leg is therefore replaced by
a cubic B-spline whose control points start on the A* path, one per
metre, and may each move up to $1.0$~m sideways. A* fixes which way
round the walls a leg goes; PSO \cite{kennedy1995} only reshapes the
curve within that corridor. A particle is the vector of sideways
offsets, and the swarm minimises a cost evaluated along the curve,
\begin{equation}
\begin{aligned}
    J = \int \Big[\, 1 &+ k_{\kappa}\kappa^{2}
      + k_{t}\,\big(|\kappa| - \kappa_{\max}\big)_{+}^{2} \\
      &+ k_{s}\,\big(c_{s} - c\big)_{+}^{2}
      + k_{h}\,\big(c_{h} - c\big)_{+} \Big]\, ds,
\end{aligned}
\label{eq:psocost}
\end{equation}
where $s$ is arc length, so the first term is the path length, $\kappa$
is the curvature, $c$ is the distance to the nearest wall read from the
clearance field, and $(\cdot)_{+} = \max(0,\cdot)$. The second term
favours gentle curves, the third penalises any turn tighter than the
minimum radius, $\kappa_{\max} = 1/r_{\min}$ with $r_{\min} = 1.0$~m,
and the last two penalise coming within a soft clearance
$c_{s} = 0.90$~m of a wall and, far more heavily, within the hard
clearance $c_{h} = 0.60$~m. One particle starts with zero
offsets, so the A* path is always among the candidates, and a leg whose
best curve still breaches $c_{h}$ falls back to its A* path. Smoothing
can therefore never make a leg less safe than the route it started from.

\begin{figure}[!t]
    \centering
    \reportfig{fig_plan_k3.png}
    \caption{Planned mission: raw A* route and the planned path after PSO smoothing. Waypoints are
    numbered in visiting order.}
    \label{fig:plan}
\end{figure}

\subsection{Parameters}
The swarm uses $30$ particles for $150$ iterations per leg, with
inertia falling from $0.9$ to $0.4$ and acceleration weights of $1.5$,
and the cost weights are $k_{\kappa} = 1$, $k_{t} = 100$, $k_{s} = 20$
and $k_{h} = 10^{4}$, so that a breach of the hard clearance outweighs
every other term. The curvature limit follows from the controller limits of Section~\ref{sec:guidance}:
the tightest turn the vehicle can hold is the speed limit over the
turn-rate limit, $v_{\max}/\omega_{\max} = 0.80$~m, and $r_{\min} = 1.0$~m
leaves turn rate in reserve to correct tracking error during a turn.

\subsection{Advantages and Disadvantages}
A* is complete and optimal on the grid: if a route exists at the
required clearance it is found, and the same waypoints always give the
same route. Its weaknesses are that the path is only optimal to within
the eight grid directions, and that its cost grows with the number of
grid points, about $2.1 \times 10^{5}$ here. Building the distance
matrix took about $15$~s of the $17$~s planning time on a laptop.

PSO removes the first weakness. It needs no gradient, so the clearance
term is read straight from a lookup table, and it penalises turns
tighter than the vehicle can make: the sharpest point of each mission's
path had a curvature of $0.78$~m$^{-1}$ on average, and one leg of $125$
exceeded the limit of $1.0$~m$^{-1}$, by $1.4$\%. It is stochastic and gives no guarantee of a global
optimum, but both are contained by design: the $\pm 1.0$~m band keeps
every candidate in the corridor A* chose, and the fallback keeps the A*
leg when smoothing fails. Over the $125$ legs of
Section~\ref{sec:results} the fallback was never needed, and smoothing
shortened the path by $1.3$\% on average.

\subsection{Reasons for Choosing the Algorithm}
The walls are known and the map is small, so an exact grid search is
affordable and its guarantees are worth having. Dijkstra's algorithm
returns the same paths as A* but searches in every direction: on the
$40$ legs of missions $1$ to $8$ it expanded $3.7$ times as many grid
points for identical path lengths, which matters when fifteen pairs must
be routed. RRT and its optimal
variant RRT* \cite{karaman2011} are built for large or high-dimensional
spaces; in a two-dimensional arena with narrow doorways they are only
probabilistically complete, can miss a doorway with a finite number of
samples, and return a different route on every run. A* also gives the
true driving distance between any two points, which the waypoint
ordering needs. PSO was added because the corners of a raw A* path
demand a turn of zero radius, which the vehicle can only approximate by
cutting into the clearance margin.

\subsection{Waypoint Ordering}
The visiting order is found with the waypoint optimiser of Part~A
\cite{dorigo1997, partA_group}, unchanged except for its cost. In
Part~A the field was empty and the cost between two points was the
straight-line distance. Here two waypoints that are close together can be
separated by a wall, so the cost $D_{ij}$ is the length of the A* path
between points $i$ and $j$, the distance the vehicle must actually drive.
All fifteen pairs among the start and the five waypoints are routed to
build $D$, and the optimiser runs on $D$ as before. The change matters:
ordering by straight-line distance would have given a longer drive in
$13$ of the $25$ missions of Section~\ref{sec:results}, by $8.6$\% averaged
over all $25$ and up to $32$\%.
